\documentclass[a4paper,12pt,oneside]{book} % Formato documento: foglio a4, font size 12, solo fronte
\usepackage[utf8]{inputenc} % Codifica caratteri
\usepackage[lmargin=2cm,rmargin=2cm,tmargin=2.5cm,bmargin=2cm]{geometry} % Bordi pagina come da formato tesi Unibs
\linespread{1.5} % Interlinea
\setlength{\parindent}{0pt} % No spazio prima dei paragrafi
\setlength{\parskip}{1em}
\setlength{\headheight}{15pt}
\addtolength{\topmargin}{-2.5pt}

\usepackage{graphicx} % Per poter inserire immagini
\usepackage{float} % Per controllare il posizionamento di immagini e tabelle
\usepackage{amsmath} % Per inserire equazioni
\usepackage{mathtools} % Per la generazione della graffa destra
\usepackage[hidelinks]{hyperref} % Per inserire URL
\usepackage{xurl} % Per consentire agli URL di spezzarsi a capo
\usepackage{enumitem} % Per poter inserire liste
\usepackage{siunitx} % Necessario per usare \setitemize
\setitemize{noitemsep,topsep=0pt,parsep=0pt,partopsep=0pt} % Per rendere le liste più compatte
\setenumerate{noitemsep,topsep=0pt,parsep=0pt,partopsep=0pt} % Per rendere le liste più compatte
\usepackage{booktabs}
\usepackage{multirow}

\usepackage[titles]{tocloft}
\setlength{\cftbeforechapskip}{6pt}

\usepackage[backend=biber,sorting=none]{biblatex} % Per poter scrivere la bibliografia su un file a parte
\addbibresource{refs.bib} % File contentente la bibliografia

\usepackage{titlesec} % Per modificare il formato dei capitoli
\titleformat{\chapter}[hang]{\normalfont\Huge\bfseries}{\thechapter.}{8pt}{\Huge\bfseries}
\titlespacing*{\chapter}{0pt}{-40pt}{20pt}


\usepackage{fancyhdr} % Per modificare l'intestazione delle pagine
\pagestyle{fancy} % Customizzazione testata e piè di pagina
\renewcommand{\chaptermark}[1]{\markboth{#1}{#1}}
\fancyhf{}
\fancyhead[L]{\thechapter. \leftmark}
\fancyfoot[C]{\thepage}


\title{Implementation of The Multiple Knapsack Problem in the Kernel Search}
\author{Matteo Beatrice, Luca Cotti, Giacomo Golino}
\date{December 2020}

\begin{document}
    \begin{titlepage}
    \begin{figure}[H] % Logo Unibs
        \centering
        \includegraphics[width=72.4mm]{images/logo_unibs}
    \end{figure}

    \begin{center}
        \LARGE{\uppercase{Università degli Studi di Brescia}}\\ % Doppio backslash = vai a capo
        \vspace{5mm} % Spazio verticale
        \large{\uppercase{Dipartimento di Ingegneria dell'informazione}}\\
        \vspace{5mm}
        \large{Corso di Laurea Magistrale in Ingegneria Informatica}\\
    \end{center}

    \vspace{10mm}

    \begin{center}
        \LARGE{\textbf{Implementation of The Multiple Knapsack Problem in the Kernel Search}}\\
    \end{center}

    \vspace{10mm}

    \begin{flushleft}
        \large
        \textbf{Docente:}\\
        Ch.ma Prof.sa Renata Mansini
    \end{flushleft}

    \begin{flushright}
        \large
        \textbf{Studenti:}\\
        Matteo Beatrice (739848)\\
        Luca Cotti (719204)\\
        Giacomo Golino (719210)
    \end{flushright}

    \vspace*{\fill} % Spazio bianco per mettere il testo in fondo alla pagina

    \rule{0.8\textwidth}{0.6pt}\\ % Linea nera
    \centering{\Large{Anno Accademico 2021/2022}}
\end{titlepage} %Frontespizio
    \hyphenpenalty=2000 % Rende meno probabile lo spezzamento delle parole a capo

    \frontmatter
    \pagenumbering{Roman} % Numeri romani maiuscoli per numero pagina per indice e intro
    \tableofcontents % Indice
    \chapter{Introduction}

In this report we describe the process of adapting the Kernel Search,
a generic heuristic framework that can be successfully used on a variety of
problems, to the specific model of the Multiple Knapsack Problem.\\
While Kernel Search provides good solutions in a short amount of time,
tuning for a specifi problem becomes necessary to improve the performance of the algorithm
and the quality of the solution found.

In the first chapter, the Multiple Knapsack Problem is introduced,
along with a few real life examples to better explain the model.\\
Chapter~\ref{ch:kernel-search} contains a brief description of the Kernel Search algorithm,
along with a few basic improvements to the default algorithm.\\
In chapter~\ref{ch:implementation} we describe the technical
specifications of our implementation of the Kernel Search algorithm,
while in chapter~\ref{ch:computational-experiments} we study
the performance of the implementation and the test instances used as benchmark.\\
In chapter~\ref{ch:improvements} we detail all
the methods that we have introduced to better adapt the Kernel Search
algorithm to the Multiple Knapsack Problem.\\
Finally, in chapter~\ref{ch:testing} we report the results of the testing we have
executed on the modified Kernel Search algorithm.
 % Introduzione

    \mainmatter
    \chapter{The Multiple Knapsack Problem}
Given a set of \(m\) knapsacks with a known \(\text{capacity } c_{i}\) \((i=1,\dots,m)\) and a set of \(n\) items with
known \(\text{profit } p_{j}\) and \(\text{weight } w_{j}\) \((j=1,\dots,n)\) the Multiple Knapsack Problem (MKP) consists in
selecting \(m\) disjoint subsets of items (one subset per knapsack) such that the total weight of the items in the knapsack
does not exceed its capacity, and the profit of the selected items is maximized.

\section{Mathematical model}
There are quite a few formulations for the MKP, with different degrees of
complexity and performance.
The formulation that was used in this project is the classical and most intuitive
one, that uses the following binary variables:
\[x_{ij}=
\begin{cases}
    1 & \quad \text{if item } j \text{ is packed into knapsack } i\\
    0 & \quad \text{otherwise}
\end{cases}
\]

The model is as follows:
\begin{equation}
    \label{eq:obj}
    \max{z} = \sum_{i=1}^{n} \sum_{j=1}^{n} p_{j} x_{ij}
\end{equation}
\begin{equation}
    \label{eq:cap}
    \sum_{j=1}^{n} w_{j} x_{ij} \leq c_{i} \qquad i=1,\dots,m
\end{equation}
\begin{equation}
    \label{eq:sel}
    \sum_{i=1}^{m} x_{ij} \leq 1 \qquad j=1,\dots,n
\end{equation}
\begin{equation}
    x_{ij} \in \{0,1\} \qquad i=1,\dots,m \quad j=1,\dots,n
\end{equation}

The objective function~\eqref{eq:obj} maximises the profit of the items in the knapsacks.\\
Constraints~\eqref{eq:cap} impose that the capacity of each knapsack is respected, while constraints~\eqref{eq:sel}
ensure that each item is packed in no more than one knapsack.

\section{Real life examples}
The MKP has various real life examples, in various fields.

The most straightforward example might be the cargo organization for a transport company.
In this scenario, the knapsack may be the trucks or containers.
Their capacity is the effective capacity of the truck/container.
The items are the objects that must be delivered.
Their weight is the actual weight of the item, while the profit is the priority of the delivery.

Another example might be the project selection:
The knapsacks may be the years of development of the project, and their capacity is the available budget for
that year.
The items can be the available projects: the weight is the cost of the project, the profit is the monetary gains
from the project completion.
    \chapter{Kernel Search}\label{ch:kernel-search}
Kernel Search is a heuristic framework used to solve MIP problems that was introduced in~\cite{kernel:2007}.\\
The algorithm has been studied in a few applications, such as the \textit{Multidimensional Knapsack Problem}
in~\cite{kernel:2010} and the \textit{Portfolio Selection Problem} in~\cite{kernel:2012}.

The algorithm is based on a few observations that can often be made when solving MIP problems:
\begin{itemize}
    \item In the optimal solution there are only a few non-zero variables;
    \item Basic variables in the LP-relaxation are good predictors of non-zero variables in the MIP optimum;
    \item Reduced costs are good predictors of the likelihood of a non-basic variable to be in the MIP optimum.
\end{itemize}

Kernel search needs an MIP solver such as GUROBI or CPLEX to solve a set of MIP sub-problems.


\section{Algorithm}
The algorithm has two main phases: an \textit{inizialitation phase}
that builds an initial \textit{kernel set} and a number of \textit{buckets},
and an \textit{improvement phase} that iteratively enlarges the kernel set and improves the solution.

\subsection{Inizialization Phase}
The first step of the Kernel Search is to solve the LP relaxation of the problem.\\
Then, the set of variables of the problem is sorted according to some \textit{sorting criterion}.

Once this is done, the kernel set \(\Lambda\) is built by selecting the \textit{C}
variables in the sorted set according a \textit{kernel construction criterion},
where \textit{C} is the size of the kernel set.

The next step is to solve MIP(\(\Lambda\)),
which is the original problem restricted to only include the variables in \(\Lambda\):
all the variables not in \(\Lambda\) are set to 0 to exclude them from the problem.\\
The resulting solution will be called \(x^{*}\) and the optimum value \(w^{*}\).\\
It is possible that no solution is found for MIP(\(\Lambda)\): in this case \(w^{*}\) is set to \(-\inf\).

The variables not in \(\Lambda\) are partitioned into a certain number of buckets \textit{nb}, according
to a \textit{bucket construction criterion}.

\subsection{Improvement Phase}
The main objectives of this phase are finding a new improving feasible solution and identifying new variables
to enter the kernel set.

Once the buckets are constructed, the algorithms proceeds with the improvement phase:\\
for each bucket \(B_{i} \ (i=1,\dots,nb)\),
the restricted sub-problem MIP(\(\Lambda \cup B_{i}\)) is solved.

Only the feasible sub-problems with an incumbent solution \(\bar{x^{*}} \geq w^{*}\) are of interest.\\
For such sub-problems the set \(\bar{\Lambda}_{i} \coloneqq \{\text{selected variables in } B_{i}\}\) is built.\\
Then the kernel set is updated: \(\Lambda \coloneqq \Lambda \cup \bar{\Lambda}_{i}\).
Since variable are only added to the kernel, and never removed, its size increases monotonically.

After all the \textit{nb} have been iterated, the algorithm ends.


\section{Parameters}
The kernel search has quite a few configuration parameters that can be used to tune the algorithm for a specific
problem:
\begin{itemize}
    \item \textit{Variable sorting criterion}: used to sort the variables in the initialization phase.
    Ideally, if the sorting criterion is selected appropriately,
    all the significant variables for the problem should be located in \(\Lambda\) and the very first buckets.
    \item \textit{Kernel size C}: the number of items initially selected to construct the kernel.
    \item \textit{Kernel construction criterion}: used to build the kernel.
    The criterion defined in the basic kernel search consists in selecting
    the first C variables in the sorted set, but for certain problems more complex criteria can be used.
    \item \textit{Bucket construction criterion}: how the buckets are constructed.
    Also defines the number and size of the buckets.
\end{itemize}


\section{Improving Efficiency of the Basic Kernel Search}\label{sec:improving-efficiency}
To improve the performance of the basic kernel search, two constraints can be added
when solving each MIP(\(\Lambda \cup B_{i}\)):
\begin{equation}
    \label{eq:impr1}
    \text{value of the objective function} \geq w^{*}
\end{equation}
\begin{equation}
    \label{eq:impr2}
    \sum_{j \in B_{i}} x_{j} \geq 1
\end{equation}
The cutoff constraint~\eqref{eq:impr1} allows only solutions that improve on the current
objective value.
Constraints~\eqref{eq:impr2} ensure that at least one item must be selected from the bucket \(B_{i}\):
such an item will then be included in the new kernel set.


\section{Iterative Kernel Search}\label{sec:iter}
Iterative kernel search is a variation of the basic kernel search where buckets are scrolled more than once.
The algorithm is as follows:
\begin{enumerate}
    \item Execute the basic kernel search;
    \item Set \(nb \coloneqq q-1\), where \(q \coloneqq \max\{i:\bar{\Lambda}_{i} \neq \emptyset\}\);
    \item Execute the improvement phase of the algorithm;
    \item If a new variable enters the kernel set, \(nb\) is reset to its initial value,
    the improvement phase is executed again and the algorithm is repeated from step 2.
\end{enumerate}

A simpler version of this algorithm repeats step 3 a certain number of times,
skipping steps 2 and 4 entirely.
The \textit{number of iterations} is a configuration parameter.
    \chapter{Implementation and Test Instances}


\section{Kernel Search}
The code provided during the course implements a generic kernel search
which uses GUROBI as solver, and has the following default paramenters:
\begin{itemize}
    \item Presolve = 2 (aggressive presolve)
    \item Time limit = 30 seconds
    \item Threads = 12
    \item MIPGap = \num{1e-12}
    \item Sort variables by profit and reduced cost
    \item Kernel size = 15\% of the number of variables (rounded to the nearest integer)
    \item Bucket size = 2.5\% of the number of variables (rounded to the nearest integer)
    \item Number of iterations = 2
    \item Time limit for each bucket = 5 seconds
\end{itemize}
The code was initially modified to implement the MKP model, and it was then
refactored to improve its quality and efficiency.
The algorithm itself and the parameters were not changed at this time.\\
The code along with extensive documentation can be found at
\url{https://github.com/Golino98/KernelSearchGolinoCottiBeatrice}.

\section{Test Instances}
The test instances for the MKP were kindly provided to us by the authors of~\cite{mkp:2019}.\\
There are in total 2100 test instances, organized in
five directories of increasing complexity:
\begin{itemize}
    \item SMALL
    \item FK\_1
    \item FK\_2
    \item FK\_3
    \item FK\_4
\end{itemize}
SMALL contains 180 instances, while FK\_1, FK\_2, FK\_3, FK\_4 contain 480 instances each.
The instances in the first three directories were already solved to the optimum.
Of those in the last two folders only a few were already solved: the unsolved
ones only had an upper and lower bound.
\newpage
\section{How to read an Instance}
Every instance was generated randomly by the authors of~\cite{mkp:2019}.
The name of each instance contains different information, in order to fastly understand the dimension of the problem represented.
The first integer represents the knapsacks number, while the second indicates then number of the items.
The last integer is used to differentiate the instances belonging to the same class.
Our test instances are text files, in which are contained the data related to the knapsacks and their \textbf{capacity} and items with their \textbf{profit} and \textbf{weight}. The structure of each of them shown below:

	\begin{flushleft} 
		$\begin{rcases*}
			3 \quad
		\end{rcases*} \rightarrow \textit{\textbf{m} knapsacks.}$
	\end{flushleft}
	\begin{flushleft}
		$\begin{rcases*}
			7 \quad
		\end{rcases*} \rightarrow \textit{\textbf{n} items.}$
	\end{flushleft}

	\begin{flushleft}
		$\begin{rcases*}
			908 \quad
			\\834 
			\\675 
		\end{rcases*} \rightarrow \textit{\textbf{capacities} of the m knapsacks.}$
	\end{flushleft}
	\begin{flushleft}
		$\begin{rcases*}
			264	\quad 430 \quad
			\\606 \quad	945
			\\268 \quad	409
			\\619 \quad	591
			\\958 \quad	839
			\\972 \quad	818
			\\723 \quad	71
		\end{rcases*} \rightarrow \textit{\textbf{weight} and \textbf{profit} of the n items.}$
	\end{flushleft}

For most of the instances without the exact solution, we were able to
find the exact using an exact GUROBI solver, with the following configuration:
\newpage
\begin{itemize}
    \item Presolve = 2 (aggressive presolve)
    \item Time limit = 90 hours
    \item Threads = 12
    \item MIPGap = \num{1e-12}
\end{itemize}

The code can be found at \url{https://github.com/Golino98/ProgettoTestingPorzioniCodiceKernelSearch}

To simplify the testing of the changes to the kernel search algorithm,
we only considered 40 instances: 6 for each directory,
plus 10 that couldn't be solved to the optimum.

Tables~\ref{tab:exact_opt} and~\ref{tab:exact_noopt} represent,
respectively, the optimum values of the 30 instances solved to
the optimum, and the upper and lower bound for the 10
that couldn't be solved.

\begin{table}[]
    \centering
    \resizebox{0.6\columnwidth}{!}{
        \begin{tabular}{@{}lll@{}}
            \toprule
            Directory              & Instance                                   & OPT    \\ \midrule
            \multirow{6}{*}{FK\_1} & random10\_60\_1\_1000\_1\_12               & 23064  \\
            & random10\_100\_2\_1000\_1\_10              & 29800  \\
            & random12\_48\_3\_1000\_1\_16               & 11865  \\
            & random15\_45\_1\_1000\_1\_13               & 15160  \\
            & random15\_75\_3\_1000\_1\_16               & 18321  \\
            & random30\_60\_4\_1000\_1\_14               & 11017  \\ \midrule
            \multirow{6}{*}{FK\_2} & random20\_120\_1\_1000\_1\_12              & 47823  \\
            & random20\_200\_2\_1000\_1\_14              & 53618  \\
            & random24\_96\_3\_1000\_1\_18               & 23912  \\
            & random30\_90\_4\_1000\_1\_20               & 26038  \\
            & random30\_150\_2\_1000\_1\_17              & 44518  \\
            & random60\_120\_4\_1000\_1\_19              & 20463  \\ \midrule
            \multirow{6}{*}{FK\_3} & random30\_180\_1\_1000\_1\_20              & 75618  \\
            & random30\_300\_2\_1000\_1\_4               & 85147  \\
            & random36\_144\_3\_1000\_1\_17              & 39245  \\
            & random45\_135\_4\_1000\_1\_16              & 37005  \\
            & random45\_225\_3\_1000\_1\_15              & 56199  \\
            & random90\_180\_4\_1000\_1\_19              & 30047  \\ \midrule
            \multirow{6}{*}{FK\_4} & random50\_300\_1\_1000\_1\_16              & 119973 \\
            & random50\_500\_2\_1000\_1\_14              & 135583 \\
            & random60\_240\_3\_1000\_1\_14              & 61180  \\
            & random75\_225\_4\_1000\_1\_17              & 60458  \\
            & random75\_375\_2\_1000\_1\_5               & 101917 \\
            & random150\_300\_2\_1000\_1\_12             & 55707  \\ \midrule
            \multirow{6}{*}{SMALL} & probT1\_0U\_R50\_T002\_M010\_N0040\_seed06 & 16510  \\
            & probT1\_0U\_R50\_T002\_M020\_N0060\_seed01 & 26593  \\
            & probT1\_0U\_R50\_T002\_M010\_N0040\_seed05 & 15534  \\
            & probT1\_1W\_R50\_T002\_M010\_N0040\_seed09 & 12724  \\
            & probT1\_1W\_R50\_T002\_M010\_N0060\_seed10 & 19652  \\
            & probT1\_1W\_R50\_T002\_M020\_N0020\_seed10 & 3429   \\
            & probT1\_1W\_R50\_T002\_M020\_N0040\_seed10 & 12405  \\ \bottomrule
        \end{tabular}
    }
    \caption{Exact solutions for the 30 instances already solved to the optimum}
    \label{tab:exact_opt}
\end{table}

\begin{table}[!htbp]
    \centering
    \begin{tabular}{@{}llll@{}}
        \toprule
        Directory              & Instance                      & LB    & UB    \\ \midrule
        \multirow{5}{*}{FK\_3} & random30\_180\_2\_1000\_1\_2  & 47911 & 47966 \\
        & random36\_144\_2\_1000\_1\_20 & 42988 & 43004 \\
        & random45\_135\_3\_1000\_1\_12 & 37680 & 37681 \\
        & random45\_135\_4\_1000\_1\_14 & 52666 & 33756 \\
        & random45\_225\_2\_1000\_1\_20 & 68070 & 68125 \\ \midrule
        \multirow{5}{*}{FK\_4} & random50\_300\_2\_1000\_1\_5  & 80459 & 80483 \\
        & random60\_240\_1\_1000\_1\_17 & 95602 & 95792 \\
        & random60\_240\_2\_1000\_1\_15 & 66497 & 66529 \\
        & random60\_240\_3\_1000\_1\_18 & 59637 & 59648 \\
        & random75\_225\_4\_1000\_1\_20 & 58338 & 58340 \\ \bottomrule
    \end{tabular}
    \caption{Upper and lower bounds for the 10 instances not solved to the optimum}
    \label{tab:exact_noopt}
\end{table}



\section{Performance of the default kernel search}


\section{Performance of random variable sorting}

    \chapter{Computational Experiments}\label{ch:computational-experiments}
In this section we report the results of tests executed on the default Kernel Search applied to the MKP\@.
The main objective of these experiments is to determine the performance of the default Kernel Search,
both in terms of quality of the solution found and time required for the execution.


\section{Test Instances}
The test instances for the MKP were kindly provided to us by the authors of~\cite{mkp:2019}.\\
There are in total 2100 test instances, organized in
five directories of increasing complexity:
\begin{itemize}
    \item SMALL;
    \item FK\_1;
    \item FK\_2;
    \item FK\_3;
    \item FK\_4.
\end{itemize}
SMALL contains 180 instances, while FK\_1, FK\_2, FK\_3, FK\_4 contain 480 instances each.
The instances in the first two directories were already solved to the optimum.
Of those in the last three folders only a few were already solved: the unsolved
ones only had an upper and lower bound.

\subsection{Format of the Instances}
The instances are \texttt{.txt} files whose name of each instance contains different information,
which allows to rapidly understand the parameters of the instance represented.

An example of the name of an instance is the following:
\begin{center}
    \textit{random10\_60\_1\_1000\_1\_12}
\end{center}
The first word is used to indicate if the instance's values are generated randomly or if they are variables of a real
life problem.
The first integer represents the knapsacks number, while the second indicates then number of the items.
The last integer is used to differentiate the instances belonging to the same class.
This instance has 10 knapsacks and 60 items, and it's the \(12^{th}\) in its class instance.

An example of the structure of an instance is shown below:
\[
    \begin{split}
        &\begin{rcases*}
             3 \quad
        \end{rcases*}
        \rightarrow \text{integer that indicates the value of \textit{m (number of knapsacks)}.}\\
        &\begin{rcases*}
             7 \quad
        \end{rcases*}
        \rightarrow \text{integer that indicates the value of \textit{n (number of items)}.}\\
        &\begin{rcases*}
             908 \quad\\
             834\\
             675\\
        \end{rcases*}
        \rightarrow \text{integers that indicate the capacity } c_{i} \text{ of the knapsack i with }
        \space i=1,\dots,m.\\
        &\begin{rcases*}
             264    \quad 430 \quad\\
             606 \quad    945\\
             268 \quad    409\\
             619 \quad    591\\
             958 \quad    839\\
             972 \quad    818\\
             723 \quad    71\\
        \end{rcases*}
        \rightarrow
        \begin{rcases*}
            \text{First column shows the weight} w_{i} \quad \\
            \text{Second column shows the profit} p_{i}.
        \end{rcases*}
        \rightarrow
        \space i=1,\dots,n.
    \end{split}
\]

The execution times required to solve the instances to the optimum were not provided to us.
For a few of the instances we were able to find the exact solution
using a GUROBI solver with the following configuration:

\begin{itemize}
    \item \textit{Presolve}: 2 (aggressive presolve)
    \item \textit{Time limit}: 1 hour
    \item \textit{Threads}: 12
    \item \textit{MIPGap}: \num{1e-12}
\end{itemize}

Most of the instances however used all the available time without reaching the optimum.
The code and the results can be found at \url{https://github.com/Golino98/EsattoMKP}.

\subsection{Selected Instances}\label{subsec:inst}
To simplify the testing of the changes to the Kernel Search algorithm,
we only considered 40 instances: 6 solved to the optimum for each directory,
plus 10 that were not solved to the optimum.

Tables~\ref{tab:exact_opt} and~\ref{tab:exact_noopt} represent,
respectively, the optimum values of the 30 instances solved to
the optimum, and the upper and lower bound for the 10
that couldn't be solved.

\begin{table}[]
    \centering
    \resizebox{0.6\columnwidth}{!}{
        \begin{tabular}{@{}lll@{}}
            \toprule
            Directory              & Instance                                   & OPT    \\ \midrule
            \multirow{6}{*}{FK\_1} & random10\_60\_1\_1000\_1\_12               & 23064  \\
            & random10\_100\_2\_1000\_1\_10              & 29800  \\
            & random12\_48\_3\_1000\_1\_16               & 11865  \\
            & random15\_45\_1\_1000\_1\_13               & 15160  \\
            & random15\_75\_3\_1000\_1\_16               & 18321  \\
            & random30\_60\_4\_1000\_1\_14               & 11017  \\ \midrule
            \multirow{6}{*}{FK\_2} & random20\_120\_1\_1000\_1\_12              & 47823  \\
            & random20\_200\_2\_1000\_1\_14              & 53618  \\
            & random24\_96\_3\_1000\_1\_18               & 23912  \\
            & random30\_90\_4\_1000\_1\_20               & 26038  \\
            & random30\_150\_2\_1000\_1\_17              & 44518  \\
            & random60\_120\_4\_1000\_1\_19              & 20463  \\ \midrule
            \multirow{6}{*}{FK\_3} & random30\_180\_1\_1000\_1\_20              & 75618  \\
            & random30\_300\_2\_1000\_1\_4               & 85147  \\
            & random36\_144\_3\_1000\_1\_17              & 39245  \\
            & random45\_135\_4\_1000\_1\_16              & 37005  \\
            & random45\_225\_3\_1000\_1\_15              & 56199  \\
            & random90\_180\_4\_1000\_1\_19              & 30047  \\ \midrule
            \multirow{6}{*}{FK\_4} & random50\_300\_1\_1000\_1\_16              & 119973 \\
            & random50\_500\_2\_1000\_1\_14              & 135583 \\
            & random60\_240\_3\_1000\_1\_14              & 61180  \\
            & random75\_225\_4\_1000\_1\_17              & 60458  \\
            & random75\_375\_2\_1000\_1\_5               & 101917 \\
            & random150\_300\_2\_1000\_1\_12             & 55707  \\ \midrule
            \multirow{6}{*}{SMALL} & probT1\_0U\_R50\_T002\_M010\_N0040\_seed06 & 16510  \\
            & probT1\_0U\_R50\_T002\_M020\_N0060\_seed01 & 26593  \\
            & probT1\_0U\_R50\_T002\_M010\_N0040\_seed05 & 15534  \\
            & probT1\_1W\_R50\_T002\_M010\_N0040\_seed09 & 12724  \\
            & probT1\_1W\_R50\_T002\_M010\_N0060\_seed10 & 19652  \\
            & probT1\_1W\_R50\_T002\_M020\_N0020\_seed10 & 3429   \\
            & probT1\_1W\_R50\_T002\_M020\_N0040\_seed10 & 12405  \\ \bottomrule
        \end{tabular}
    }
    \caption{Exact solutions for the 30 instances already solved to the optimum}
    \label{tab:exact_opt}
\end{table}

\begin{table}[!htbp]
    \centering
    \begin{tabular}{@{}llll@{}}
        \toprule
        Directory              & Instance                      & LB    & UB    \\ \midrule
        \multirow{5}{*}{FK\_3} & random30\_180\_2\_1000\_1\_2  & 47911 & 47966 \\
        & random36\_144\_2\_1000\_1\_20 & 42988 & 43004 \\
        & random45\_135\_3\_1000\_1\_12 & 37680 & 37681 \\
        & random45\_135\_4\_1000\_1\_14 & 52666 & 33756 \\
        & random45\_225\_2\_1000\_1\_20 & 68070 & 68125 \\ \midrule
        \multirow{5}{*}{FK\_4} & random50\_300\_2\_1000\_1\_5  & 80459 & 80483 \\
        & random60\_240\_1\_1000\_1\_17 & 95602 & 95792 \\
        & random60\_240\_2\_1000\_1\_15 & 66497 & 66529 \\
        & random60\_240\_3\_1000\_1\_18 & 59637 & 59648 \\
        & random75\_225\_4\_1000\_1\_20 & 58338 & 58340 \\ \bottomrule
    \end{tabular}
    \caption{Upper and lower bounds for the 10 instances not solved to the optimum}
    \label{tab:exact_noopt}
\end{table}


\section{Tests of the default Kernel Search}
The following table outlines the result of the Kernel Search
algorithm with the default configuration.
The tables includes the best solution found, the time elapsed in order to find it (with a maximum time limit of two minutes) and if the time limit was reached.

\begin{table}[!htbp]
    \centering
    \resizebox{0.7\columnwidth}{!}{
        \begin{tabular}{@{}lllll@{}}
            \toprule
            Directory              & Instance                                   & OPT     & Time Elapsed    &Is the optimum   \\ \midrule
            \multirow{6}{*}{SMALL}
            & probT1\_0U\_R50\_T002\_M010\_N0040\_seed05 & 15378 & 7.7404257 & false \\  
            & probT1\_0U\_R50\_T002\_M020\_N0060\_seed01 & 25970 & 20.687005 & false \\  
            & probT1\_1W\_R50\_T002\_M010\_N0040\_seed09 & 12664 & 14.5433935 & false \\  
            & probT1\_1W\_R50\_T002\_M010\_N0060\_seed10 & 19589 & 119.6407061 & true \\  
            & probT1\_1W\_R50\_T002\_M020\_N0020\_seed10 & 3429 & 4.763744 & false \\  
            & probT1\_1W\_R50\_T002\_M020\_N0040\_seed10 & 11951 & 10.6224961 & false \\   
            \midrule
            \multirow{6}{*}{FK\_1} 
            & random10\_100\_2\_1000\_1\_10 & 29698 & 118.1525164 & true \\  
            & random10\_60\_1\_1000\_1\_12 & 23064 & 18.3776223 & false \\  
            & random12\_48\_3\_1000\_1\_16 & 11813 & 52.980434 & false \\  
            & random15\_45\_1\_1000\_1\_13 & 14926 & 11.1281064 & false \\  
            & random15\_75\_3\_1000\_1\_16 & 18298 & 118.1294058 & true \\  
            & random30\_60\_4\_1000\_1\_14 & 11017 & 29.2274437 & false \\  
            \midrule
            \multirow{6}{*}{FK\_2} 
            & random20\_120\_1\_1000\_1\_12 & 47764 & 120.2145511 & true \\  
            & random20\_200\_2\_1000\_1\_14 & 53472 & 118.1115582 & true \\  
            & random24\_96\_3\_1000\_1\_18 & 23861 & 118.2848323 & true \\  
            & random30\_150\_2\_1000\_1\_17 & 44240 & 118.0832715 & true \\  
            & random30\_90\_4\_1000\_1\_20 & 25977 & 55.9372853 & false \\  
            & random60\_120\_4\_1000\_1\_19 & 20463 & 118.1919825 & true \\  
            \midrule
            \multirow{5}{*}{FK\_3\textbackslash notSolved} 
            & random30\_180\_2\_1000\_1\_2 & 47752 & 119.5763461 & true \\  
            & random36\_144\_2\_1000\_1\_20 & 42583 & 118.4716056 & true \\  
            & random45\_135\_3\_1000\_1\_12 & 37575 & 118.2264545 & true \\  
            & random45\_135\_4\_1000\_1\_14 & 33665 & 118.0438898 & true \\  
            & random45\_225\_2\_1000\_1\_20 & 67784 & 118.2544371 & true \\  
            \midrule
            \multirow{6}{*}{FK\_3\textbackslash Solved}
            & random30\_180\_1\_1000\_1\_20 & 75456 & 118.6970463 & true \\  
            & random30\_300\_2\_1000\_1\_4 & 84846 & 120.6542013 & true \\  
            & random36\_144\_3\_1000\_1\_17 & 39165 & 118.7504888 & true \\  
            & random45\_135\_4\_1000\_1\_16 & 36953 & 118.079937 & true \\  
            & random45\_225\_3\_1000\_1\_15 & 56007 & 121.0977372 & true \\  
            & random90\_180\_4\_1000\_1\_19 & 30047 & 118.0271635 & true \\  
            \midrule
            \multirow{5}{*}{FK\_4\textbackslash notSolved}
            & random50\_300\_2\_1000\_1\_5 & 80094 & 120.648261 & true \\  
            & random60\_240\_1\_1000\_1\_17 & 95285 & 119.4535949 & true \\  
            & random60\_240\_2\_1000\_1\_15 & 66103 & 118.140116 & true \\  
            & random60\_240\_3\_1000\_1\_18 & 59406 & 118.8488718 & true \\  
            & random75\_225\_4\_1000\_1\_20 & 58306 & 118.7606362 & true \\   
            \midrule
            \multirow{6}{*}{FK\_4\textbackslash Solved}
            & random150\_300\_2\_1000\_1\_12 & 54611 & 119.06234 & true \\  
            & random50\_300\_1\_1000\_1\_16 & 119452 & 122.0148949 & true \\  
            & random50\_500\_2\_1000\_1\_14 & 134877 & 122.2307049 & true \\  
            & random60\_240\_3\_1000\_1\_14 & 60900 & 120.9684616 & true \\  
            & random75\_225\_4\_1000\_1\_17 & 60411 & 118.9970356 & true \\  
            & random75\_375\_2\_1000\_1\_5 & 101012 & 122.4775278 & true \\  
            \bottomrule
        \end{tabular}
        }
    \caption{Results of Kernel builder with integer variables sorted by value and absolute RC}
    \label{tab:ker_int_val_abs_RC}
\end{table}

    \chapter{Improvements}\label{ch:improvements}
In this section we will describe the various changes
that we applied to the default kernel search implementation,
with the objective of improving its performance
in solving the MKP instances.

For the sake of completeness, we well also describe the all attemps at improving
the algorithm, thus including the unsuccesful ones.

The effectiveness of the tuning is demonstrated using the selected instances
described in~\ref{subsec:inst} as benchmark.


%\section{Variable Sorting}


\section{Kernel Construction Criteria}
A recurring problem we have identified when running more complex
instances (such as the one in \texttt{FK\_4}) is that,
from the very beginning, the kernel contains quite a lot of variables,
which makes the solving of the kernel problem and the buckets sub-problems slow.

This means that a starting point to improve the performance of the Kernel Search
could be to optimize the management of the kernel set,
possibily starting from the \textit{kernel construction criterion},
which by default is to include the first \textit{C} sorted variables.

We have experimented with the following criteria:
\begin{enumerate}
    \item Select the variables that have a greater than 0 in the LP relaxation;
    \item Select the variables that have a value equal to 1 in the LP relaxation;
    \item Select the variables that have a value greater than a threshold in the LP relaxation.
\end{enumerate}

However, after running benchmark to compare these criteria
with the default one, the results demonstrated that the default criterion
consistently gave better results.
It is however worth noting that the criterion 2 and 3 (with a proper threshold),
gave results that were quite close to the ones with the default criteria,
while using much less time.

%\section{Ejection of Variables from the Kernel}


\section{Repetition Counter}
An observation we could make is that when solving buckets,
after a certain number of iterations the algorithm
repeatedly finds new solutions with the same objective value,
and sometimes it even struggles to find new solutions at all.
In other words, the algorithm gets stuck in \textit{local optima},
from which it's hard to escape.

Our hypothesis is that there are two causes for this:
\begin{enumerate}
    \item The GUROBI solver cannot find a (better) solution for the sub-problem,
    either due to the infeasibility of the problem or because of
    the time limit set for solving each bucket.
    \item The kernel search algorithm, during its improvement phase,
    only accepts a new solution (thus updating the kernel)
    if it improves upon, or is at least equal, to the incumbent one.
    This is done by adding a cutoff constraint, as explained in~\ref{sec:improving-efficiency}.
\end{enumerate}

In order to mitigate this problem we introduced a \textit{repetition counter}
that, during the improvement phase of kernel search,
removes the cutoff constraint for \textit{k} buckets
when the same solution (or no solution) is found for \textit{h} times.

The idea is that this method allows to select if the focus should be on
\textit{diversification} or \textit{intensification},
by appropriately setting parameters \textit{h} and \textit{k}.

A low value for \textit{h} and a high one for \textit{k} allow for
diversification, by adding to the kernel variables that otherwise
may never be selected, and could allow to escape the local optimum.

On the opposite, a high value for \textit{h} and a low one for \textit{k}
switch to focus on intensification, by giving priority to finding
better solutions.

After experimenting with different values for \textit{h} and \textit{k},
we found that keeping them more or less equal allowed for
a reasonable balance between diversification and intensification.
In particular we found that \(h=3\) and \(k=3\) worked quite well
with the test instances we selected.

\subsection{Improvements}
A small enhancement we have applied to this mechanism
is to reset the counter to its initial status
whenever the algorithm finds a new solution
that is better than any other found before.

This is significant because the default method
completely ignores the value of the solutions found
during the \textit{k} cycles.

%\section{Iterative Kernel Search}


    \backmatter
    \printbibliography[
        heading=bibintoc % Inserisci bibliografia nell'indice
    ]
\end{document}